% !Mode:: "TeX:UTF-8"
% !TEX program  = xelatex

\documentclass[bwprint]{report}
%\documentclass[bwprint]{hfutreport} %打印的时候有超链接的地方不需要彩色，可以加上bwprint 选项

%\usepackage{bibentry,natbib}

\title{基于生成模型的图像超分辨率方法研究与轻量化改进设计}
\stunum{第X组}
\stuname{{成员A、成员B、成员C}}
\stuclass{计算机20XX届X班}
\supervisor{XXX \hspace{1.0em} 老师}
\dateinput{\today}
%\maincontent{针对数字图像处理领域某一特定问题，例如：智能问答、信息检索、生成式对话、信息抽取、推荐系统、知识图谱等，查阅国内外近5年来相关技术资料，选择3篇以上具有代表性论文进行深入阅读，分析并理解其原理和方法。并撰写完成一篇不低于4000字的课程设计报告。对相关论文的理论原理及算法进行阐述，对其中的方法动手搭建模型，编写程序复现实验结果，进行对比分析，给出结论。自行组队，每组1~3人。报告中明确每人的分工。考核要点包括：
%（1）逻辑清晰，主题明确，图文并茂，内容丰富；
%（2）选题新颖，阅读和介绍近年来较新的相关研究成果；
%（3）试验过程详细，图表对比分析合理；
%（4）算法理解分析透彻；
%（5）提出自己的算法设想，编写程序实现仿真实验；
%（6）说明哪些是别人已有的工作，哪些是自己做的工作。
%}

\begin{document}

 \maketitle

% 目录使用Romen页脚
\pagestyle{plain}
\pagenumbering{Roman}
\tableofcontents
\newpage
\pagestyle{plain}

%正文页脚
\setcounter{page}{1}
\pagenumbering{arabic}


\section{团队成员分工}
\begin{itemize}
    \item 成员A：负责文献检索、相关工作梳理、任务背景与任务描述撰写。
    \item 成员B：负责模型实现、实验设计、结果整理与图表绘制。
    \item 成员C：负责改进方案设计、总结撰写及 \LaTeX{} 排版整合。
\end{itemize}



\section{任务背景}
图像超分辨率（Super-Resolution, SR）是计算机视觉中的经典研究问题，其目标是从低分辨率图像中恢复出更高分辨率的图像细节。该任务在医学图像增强、遥感图像处理、视频增强、安防监控以及老旧图像修复等场景中具有重要应用价值。传统方法通常依赖双线性插值、双三次插值等手工设计的重建策略，但这类方法往往只能恢复平滑区域，难以有效重建高频纹理与边缘细节。

随着深度学习的发展，图像超分辨率方法逐渐从卷积神经网络模型发展到基于生成对抗网络和扩散模型的生成式方法。卷积神经网络方法强调像素级重建精度，在峰值信噪比和结构相似性等指标上具有较好表现；生成对抗网络方法则更注重视觉感知质量，能够生成更加锐利的纹理细节；近年来扩散模型进一步推动了高质量图像恢复的发展，使生成式超分辨率成为研究热点。因此，对生成模型在图像超分辨率中的方法进行系统梳理，并在轻量模型上进行复现与改进，具有较强的理论意义和实践价值。

\section{任务描述}
本课程设计围绕“基于生成模型的图像超分辨率方法研究”展开。首先，通过查阅近五年国内外代表性研究论文，分析图像超分辨率领域中轻量重建模型、生成对抗网络模型以及扩散生成模型的发展脉络与核心思想；其次，结合课程设计的实验条件，选择一个计算开销较低、结构相对清晰的轻量模型作为主要复现对象，并基于 MindSpore 框架完成模型搭建；最后，在主模型基础上引入生成模型中的感知增强思想，设计轻量改进模块，并通过定量与定性实验对方法效果进行分析。

本课题的重点不在于完整训练复杂的大规模生成模型，而在于：一是对近年来代表性超分辨率方法的原理进行深入分析；二是完成一个可复现的轻量模型实现；三是在实验资源受限的情况下提出合理的改进思路，并完成一定程度的实验验证。通过该课程设计，可以系统理解超分辨率任务的研究背景、技术路线及生成模型在视觉计算中的实际应用。

\section{技术方案}
本文采用“文献综述 + 轻量复现 + 改进设计”的技术路线。整体思路是先对超分辨率方法的发展进行归纳，再以轻量模型作为实验主线完成复现，最后借鉴生成模型中的感知增强思想进行轻量化改进。这样既能够体现课程设计对前沿方法的关注，又能兼顾本地实验设备的稳定性和实现难度。

\subsection{方案总体架构}
整体技术架构分为四个层次：第一层为任务与文献研究层，负责对图像超分辨率任务背景、研究现状以及典型方法进行梳理；第二层为基线模型实现层，主要基于 MindSpore 实现 SRCNN 等基础超分辨率模型；第三层为主模型复现与改进层，以 IMDN 一类轻量模型为核心，设计加入边缘损失或轻量注意力模块的改进方案；第四层为实验评估层，通过 PSNR、SSIM 及可视化结果对不同方法进行对比分析。

\subsection{模块1}
模块1为文献研究与方法分析模块。为了保证论文分析既具有前沿性，又具备清晰的技术演进逻辑，本文选取三类代表性方法作为重点研究对象，分别对应轻量重建、生成对抗网络超分以及扩散模型超分三条技术路线。通过对这三类方法的梳理，可以从“重建保真”“感知增强”“概率生成”三个层面理解近年来图像超分辨率方法的发展脉络。

第一类方法选取 IMDN 作为代表。IMDN 属于轻量高效的图像超分辨率网络，其核心思想是通过信息多蒸馏模块逐步提取和筛选有效特征，在保证较少参数量的同时维持较好的重建性能。与早期简单堆叠卷积层的超分模型相比，IMDN 更强调对特征信息的有效利用，适合作为课程设计中的主复现对象。一方面，该模型结构相对清晰，便于分析与实现；另一方面，其轻量特性也更符合当前实验设备对训练成本和稳定性的要求。

第二类方法选取 ESRGAN 作为代表。ESRGAN 是生成对抗网络在超分辨率任务中的典型工作之一，其主要目标不再局限于提高 PSNR 和 SSIM，而是更加注重图像的感知质量和纹理真实感。该方法通过改进网络结构和损失设计，使生成结果在视觉上更自然、更锐利。ESRGAN 的重要意义在于，它揭示了“高指标”和“高观感”并不总是完全一致，为后续生成式超分辨率研究提供了重要思路。对于本课程设计而言，ESRGAN 非常适合用来作为生成模型路线的核心分析对象。

第三类方法选取 SR3 作为代表。SR3 将扩散模型引入图像超分辨率任务，其基本思想是在条件约束下通过逐步去噪的方式生成高分辨率图像。相较于传统卷积方法和 GAN 方法，扩散模型在生成细节和整体视觉一致性方面表现出更大的潜力，也因此成为近年来视觉生成领域的重要发展方向。虽然 SR3 的训练和推理成本较高，不适合作为本课程设计的主训练模型，但它能够很好地体现生成模型在超分辨率任务中的前沿研究趋势。

综合来看，IMDN、ESRGAN 和 SR3 分别对应轻量高效重建、感知驱动生成和扩散概率生成三条典型路线。将这三类方法放在同一框架下分析，不仅能够增强报告的逻辑完整性，也有助于说明本文为什么最终选择轻量模型作为实验主体，同时借鉴生成模型中的细节增强思想进行改进设计。

\subsection{模块2}
模块2为模型实现与改进设计模块。实验部分以 SRCNN 作为基础基线模型，以 IMDN 作为主要复现对象。SRCNN 结构简单，适合用于建立基础对照；IMDN 通过信息蒸馏思想减少冗余特征，兼顾了重建效果与模型复杂度。为了体现本课题与生成模型研究方向的联系，本文计划在轻量模型基础上加入边缘损失、感知约束或轻量注意力模块，从而提升模型的细节恢复能力，并分析这种改进在课程设计场景下的可行性与效果。

\subsection{典型论文方法原理分析}
为了更系统地说明图像超分辨率方法的发展脉络，本文选取 IMDN、ESRGAN 和 SR3 三篇具有代表性的工作进行原理分析。这三篇论文分别对应轻量重建模型、生成对抗网络模型和扩散模型三条技术路线，既能够体现超分辨率研究的阶段性演进，也能够为后续实验设计与模型改进提供理论依据。

\subsubsection{IMDN 方法原理分析}
IMDN（Information Multi-Distillation Network）是一种面向轻量化图像超分辨率任务的代表性模型，其设计目标是在较低参数量和较少计算开销的条件下，尽可能提升图像重建质量。与传统通过不断堆叠卷积层来提升性能的方式不同，IMDN 更加强调对已有特征的高效提取、筛选与融合，从而在轻量模型框架下实现较好的超分辨率效果。

IMDN 的核心思想是“信息多蒸馏”。具体而言，模型在特征提取过程中不会一次性将全部特征送入后续层，而是通过蒸馏模块逐步分离出更有价值的特征信息，并将这些信息与后续提取结果进行融合。这种做法能够有效减少冗余特征传播，提高特征利用效率。与普通卷积堆叠网络相比，IMDN 能够在更小的模型规模下保持较强的表示能力，这也是其适合部署在资源受限设备上的重要原因。

从网络结构来看，IMDN 通常由浅层特征提取、多个信息蒸馏模块、特征融合模块以及上采样重建模块组成。浅层卷积负责从输入低分辨率图像中提取基础纹理信息；中间的蒸馏模块通过逐步提取和筛选特征，保留有利于图像恢复的关键信息；最后通过上采样模块将低分辨率特征映射到高分辨率空间，输出重建图像。其整体结构紧凑、逻辑清晰，因此非常适合作为课程设计中的主复现对象。

IMDN 的优点在于结构轻量、训练相对稳定、推理速度快，能够在轻量级模型中取得较有竞争力的 PSNR 和 SSIM 指标。然而，其局限也较为明显：由于模型容量受到限制，在复杂纹理恢复和高倍率放大任务中，其细节生成能力仍然弱于更大规模的生成式模型。因此，在本课程设计中，IMDN 更适合作为实验主体，而生成模型中的细节增强思想则可以作为后续改进方向引入其中。

\subsubsection{ESRGAN 方法原理分析}
ESRGAN（Enhanced Super-Resolution Generative Adversarial Networks）是图像超分辨率领域中生成对抗网络路线的代表性方法。与传统方法主要追求较高 PSNR 和 SSIM 的思路不同，ESRGAN 更关注人眼主观感知效果，希望生成在视觉上更加真实、自然且具有丰富纹理细节的高分辨率图像。因此，ESRGAN 的提出标志着超分辨率研究从“重建精度导向”逐步转向“感知质量导向”。

ESRGAN 的核心在于将生成对抗学习思想引入图像超分辨率任务。其生成器通过学习低分辨率图像到高分辨率图像的映射关系，努力恢复清晰的纹理细节；判别器则负责区分生成结果与真实高分辨率图像，从而推动生成器不断改进输出质量。在这一过程中，模型不再仅仅以像素误差最小化为目标，而是通过对抗损失和感知损失共同约束，促使生成图像在视觉上更加逼真。

在结构设计上，ESRGAN 引入了改进的残差密集块结构，用于增强特征表达能力，并在一定程度上提高训练稳定性。与此同时，该方法还使用感知损失来约束生成图像与真实图像在高层语义特征空间中的接近程度，而不是只关注像素级差异。这种多重损失共同优化的方式，使得 ESRGAN 在纹理恢复、边缘锐化以及主观视觉效果上均优于许多传统重建型方法。

ESRGAN 的优势在于能够生成更加锐利和自然的图像细节，特别是在草地、砖墙、头发、衣物纹理等复杂区域的恢复中表现突出。然而，其不足也同样值得注意。首先，生成对抗网络训练过程相对不稳定，对超参数和训练策略较为敏感；其次，为了获得更好的视觉观感，模型可能会产生与真实图像不完全一致的伪纹理，从而导致 PSNR 等客观指标不一定最优。因此，在本课程设计中，ESRGAN 更适合作为生成模型超分路线的理论分析对象，用于说明“感知质量”和“数值指标”之间的权衡关系。

\subsubsection{SR3 方法原理分析}
SR3（Super-Resolution via Repeated Refinement）是扩散模型在图像超分辨率任务中的代表性工作之一。该方法将超分辨率问题建模为条件生成问题，即在给定低分辨率图像的条件下，通过逐步去噪和逐步细化的方式生成高分辨率图像。与传统卷积模型和 GAN 模型相比，SR3 的显著特点在于其生成过程具有更强的概率建模能力，能够更自然地恢复细节并保持整体图像结构的一致性。

SR3 的基本思想来源于扩散模型。扩散模型通常先对真实图像逐步加入噪声，再训练网络学习如何反向去噪，以便从噪声样本恢复出真实图像。在超分辨率任务中，低分辨率图像被作为条件输入，模型在多步去噪过程中不断细化图像内容，最终生成高分辨率结果。这种“逐步优化”的过程使得模型在恢复复杂纹理和高频信息时具有更大的灵活性和表现空间。

从方法机制来看，SR3 强调条件信息在生成过程中的约束作用。低分辨率图像不仅提供了整体轮廓和结构基础，也为每一步去噪过程提供引导，使最终输出既保持内容一致性，又能够具备更丰富的细节表达。相比 GAN 方法，扩散模型通常在训练稳定性和生成多样性方面具有一定优势；相比传统重建模型，其在高感知质量图像恢复方面也展现出更大的潜力。

SR3 的主要优势在于生成结果自然、细节丰富，并且在高倍率图像恢复任务中往往具有更强的视觉表现力。然而，该方法的代价也比较明显：由于生成过程需要多次迭代去噪，其推理耗时和计算成本通常高于普通卷积网络和 GAN 模型；同时，训练扩散模型对算力和实验环境的要求也更高。因此，在本课程设计中，SR3 更适合作为前沿生成式方法的理论研究对象，用于说明扩散模型在图像超分辨率中的研究趋势，而不作为主要复现模型。
实验部分围绕“低风险、可逐步推进”的原则进行设计。考虑到本地笔记本在长时间高负载运行下可能存在稳定性问题，实验将优先完成环境验证、最小前向推理和小样本验证，在设备稳定的前提下再逐步补充正式训练结果。实验结果将从定量指标与视觉效果两个角度展开分析。

\subsection{数据集}
本实验拟采用公开超分辨率数据集进行验证，例如 Set5、Set14、BSD100 或 DIV2K 的部分子集。对于低分辨率图像的构造，可采用双三次下采样方式生成对应样本。若后期训练资源允许，可在 DIV2K 子集上进行训练，并在 Set5 或 Set14 数据集上进行测试；若训练资源受限，则可先以小样本子集和最小前向验证为主，保证实验流程完整可复现。

\subsection{评测指标}
超分辨率任务常用的定量评价指标包括峰值信噪比（PSNR）和结构相似性（SSIM）。其中，PSNR 用于衡量重建图像与真实高分辨率图像之间的像素误差水平，数值越高表示重建误差越小；SSIM 用于衡量图像在结构信息上的相似程度，能够更好反映人眼感知上的一致性。除定量指标外，本文还将通过典型图像区域的视觉对比来分析模型对纹理、边缘和细节的恢复能力。

\subsection{实现细节}
本实验拟基于华为 MindSpore 深度学习框架进行实现，以满足课程设计“建议基于 MindSpore” 的要求。基础实验配置可设置为单卡 GPU、放大倍率 $\times 2$ 或 $\times 4$、优化器采用 Adam、重建损失采用 L1 或 L2 损失。改进模型部分可进一步引入边缘损失或感知增强项。由于实验设备稳定性存在不确定性，当前阶段优先完成模型结构搭建、前向推理验证、配置文件设计和训练脚本骨架编写，再逐步补充训练过程与正式实验结果。

\subsection{定量实验}
定量实验的目标是比较基线模型、主模型及改进模型在公开测试集上的数值表现。实验以 PSNR 和 SSIM 作为主要评价标准，通过统一的数据预处理与评价流程，比较不同方法在重建精度上的差异。考虑到当前阶段以小样本验证和最小可复现为主，表中结果主要用于说明不同模型的相对表现以及改进方案的可行性。

\subsubsection{与其他方法比较}
本小节比较双三次插值、SRCNN、IMDN 以及改进模型在 $\times 2$ 放大任务下的性能差异。双三次插值作为传统方法，用于提供基础参考；SRCNN 作为经典卷积超分模型，可用于说明深度学习方法相对传统插值的提升；IMDN 作为轻量高效模型，在参数量和性能之间具有较好平衡；改进模型则重点考察在引入边缘约束后对重建质量的影响。通过结果对比，可以分析轻量模型在课程设计场景中的适用性，以及借鉴生成模型思想进行轻量化改进的实际价值。

\begin{table}[!htbp]
    \caption{不同超分辨率方法在测试集上的定量比较}
    \label{tab:quantitative-comparison}
    \centering
    \begin{tabular}{lcccc}
        \toprule[1.5pt]
        方法 & 参数量 / M & 放大倍率 & PSNR / dB & SSIM \\
        \midrule[1pt]
        双三次插值 & -- & $\times 2$ & 待补充 & 待补充 \\
        SRCNN & 0.057 & $\times 2$ & 待补充 & 待补充 \\
        IMDN & 0.694 & $\times 2$ & 待补充 & 待补充 \\
        IMDN + Edge Loss & 0.694 & $\times 2$ & 待补充 & 待补充 \\
        \bottomrule[1.5pt]
    \end{tabular}
\end{table}

从表\ref{tab:quantitative-comparison}的设计可以看出，本文重点并非追求大规模模型的极限性能，而是关注轻量模型在有限实验条件下的可复现性与可改进性。若后续实验结果完整，可进一步分析 SRCNN 与 IMDN 在参数规模和重建性能之间的差异，并讨论加入边缘损失后模型在 PSNR、SSIM 指标上的变化趋势。

\subsubsection{消融实验}
为了验证改进设计的有效性，本文在 IMDN 主模型基础上引入边缘损失项，并通过与原始 IMDN 的对比分析该改进对重建性能的影响。考虑到课程设计阶段的实验资源有限，本报告将消融实验聚焦于“是否加入边缘损失”这一核心问题，而不进一步扩展过多改进组件，以保证实验结论简洁明确。

\begin{table}[!htbp]
    \caption{边缘损失消融实验结果}
    \label{tab:ablation-study}
    \centering
    \begin{tabular}{lcc}
        \toprule[1.5pt]
        实验设置 & Edge Loss & PSNR / dB \\
        \midrule[1pt]
        Baseline IMDN & 否 & 待补充 \\
        IMDN + Edge Loss & 是 & 待补充 \\
        \bottomrule[1.5pt]
    \end{tabular}
\end{table}

表\ref{tab:ablation-study}用于分析边缘约束是否能够帮助模型更好地恢复图像中的结构边界信息。若加入边缘损失后模型取得更高的 PSNR 或者在视觉效果上表现出更清晰的轮廓，则可以说明该改进在轻量超分辨率模型中具有一定有效性；若数值提升有限，也可据此说明感知增强类改进在轻量模型上的收益仍受模型容量限制。

\subsection{定性实验}
定性实验主要通过图像可视化结果展示模型在纹理恢复、边缘锐化和细节生成方面的表现。与单纯的数值指标相比，定性实验更能够体现生成式超分辨率方法在感知质量上的优势。当前阶段由于实验结果图仍在整理，本文暂不在报告中展示占位图片，后续将补充低分辨率输入、双三次插值、基线模型与改进模型的可视化对比结果。

\subsubsection{案例分析}
案例分析将从图像中的典型区域入手，例如建筑纹理、边缘线条、文字区域或自然纹理等，观察不同方法在细节恢复上的差异。通常情况下，基础重建模型更注重整体误差控制，而引入感知增强思想的改进模型在局部纹理恢复上可能更具优势。通过这些具体案例，可以更直观地分析不同模型的特点与不足。

后续补充实验结果时，将重点比较边缘区域、重复纹理区域和复杂高频区域的恢复情况。例如，在建筑立面、衣物纹理和文字边缘区域中，若改进模型能够生成更清晰的结构轮廓，则可说明其在视觉感知层面优于基础模型；若改进效果主要体现在边缘平滑或局部清晰度提升上，也能够说明边缘约束在轻量超分任务中的实际作用。

\subsubsection{结果展示说明}
考虑到当前实验仍处于逐步完善阶段，定性结果部分将优先展示具有代表性的单张样例图，并从整体清晰度、局部边缘恢复和纹理自然度三个角度进行分析。与大规模完整实验相比，这种小样本展示方式更符合本课程设计“低风险、可逐步推进”的实施策略，同时也能够较真实地反映当前模型实现与改进设计的阶段性效果。

\section{总结}
本文围绕基于生成模型的图像超分辨率方法展开研究，首先梳理了近年来轻量重建模型、生成对抗网络模型和扩散模型在超分辨率任务中的发展脉络，分析了不同方法在重建精度、视觉质量和训练复杂度方面的特点；然后结合课程设计的实验条件，选择轻量模型作为主要实现对象，并在此基础上设计了引入生成模型思想的轻量化改进方案。整体方案兼顾了研究前沿性与实现可行性，适合作为课程设计的技术路线。

在相关方法分析方面，本文重点研究了 IMDN、ESRGAN 和 SR3 三类代表性方法。IMDN 通过信息多蒸馏结构实现了较好的轻量化重建效果，可作为本文实验部分的主要复现对象\cite{hui2019imdn}；ESRGAN 通过生成对抗学习显著提升了图像纹理和视觉真实感，体现了生成模型在超分辨率任务中的重要价值\cite{wang2018esrgan}；SR3 则将扩散模型引入超分辨率任务，通过逐步去噪生成高质量结果，代表了当前生成式超分辨率的发展趋势\cite{saharia2021sr3}。

在当前阶段，本文已经完成了课题选题、文献分析框架、报告结构设计以及 MindSpore 代码骨架准备，并完成了实验方案、比较指标和改进思路的整理。后续工作将在保证设备稳定的前提下，逐步完成最小前向验证、小规模实验以及定量和定性结果整理。通过本课程设计，可以进一步加深对图像超分辨率任务与生成模型视觉应用的理解，也为后续探索更高质量、更高效率的图像恢复方法奠定基础。

\bibliographystyle{plain}
\bibliography{mybib}

\end{document}
